\begin{answer}
    % Rebuttal response to Reviewer VY3D-Q2
    % Topic: Computational overhead of SSS in long-audio decoding

    \paragraph{Response to Reviewer VY3D-Q2.}
    We thank the reviewer for raising this important point.
    We agree that KV-memory savings alone are not sufficient to characterize the practicality of SSS, especially for long-audio ASR where the audio prefix is large.
    We therefore measured the actual end-to-end decoding overhead of SSS inside the full Qwen2.5-Omni-3B ASR pipeline, rather than only reporting an isolated RFFT/IRFFT operator benchmark.

    \paragraph{End-to-end timing in real ASR decoding.}
    On LibriSpeech, we compare the same decoding setup with and without SSS enabled.
    The average wall-clock time increases from 4.091s to 4.173s per sample, corresponding to only +0.081s absolute overhead or +1.99\% relative slowdown.
    The average decode time changes from 0.0677s to 0.0691s, i.e., +2.02\%.
    In per-token terms, decoding latency increases from 0.68ms/token to 0.69ms/token, an additional 0.014ms/token.
    Overall throughput changes from 1477.47 to 1448.16 tokens/s, a 1.98\% drop.

    \begin{table}[t]
    \centering
    \small
    \setlength{\tabcolsep}{4pt}
    \caption{\textbf{End-to-end overhead of SSS in real ASR decoding.} Timing is measured on LibriSpeech using Qwen2.5-Omni-3B with identical decoding settings, comparing SSS disabled vs. enabled.}
    \label{tab:rebuttal-sss-e2e-overhead}
    \begin{tabular}{lccc}
    \toprule
    \textbf{Metric} & \textbf{w/o SSS} & \textbf{w/ SSS} & \textbf{Change} \\
    \midrule
    Avg. wall time / sample & 4.091s & 4.173s & +1.99\% \\
    Avg. decode time / sample & 0.0677s & 0.0691s & +2.02\% \\
    Per-token decode latency & 0.68ms & 0.69ms & +0.014ms \\
    Overall throughput & 1477.47 tok/s & 1448.16 tok/s & -1.98\% \\
    \bottomrule
    \end{tabular}
    \end{table}

    \paragraph{Why the practical overhead is small.}
    Although the frequency-domain smoothing operation involves RFFT/IRFFT, it is applied only to the selected score/saliency computation path and is lightweight compared with the full multimodal forward pass and attention computation over long audio contexts.
    In the full ASR pipeline, the dominant cost remains model prefill and decoding, so the additional SSS computation leads to a small end-to-end slowdown of about 2\%.

    \paragraph{Planned revision.}
    We will add Table~\ref{tab:rebuttal-sss-e2e-overhead} and the above timing analysis to the efficiency section.
    We will also clarify that isolated RFFT/IRFFT timing should be interpreted as an operator stress test, while the primary practical metric is the end-to-end ASR decoding overhead with SSS enabled.
\end{answer}
